\origpage{177}{GRAPH MINORS, X}
\stmt
(8.5) \textit{Let $\mathscr{T}$ be a tangle in a graph $G$ of order $\theta \ge 1$. Let $W \subseteq V(G)$ with $|W| < \theta$. Let $\mathscr{T}'$ be the set of all separations $(A', B')$ of $G \setminus W$ of order $< \theta - |W|$ such that there exists $(A, B) \in \mathscr{T}$ with $W \subseteq V(A \cap B)$ and $A \setminus W = A'$, $B \setminus W = B'$. Then $\mathscr{T}'$ is a tangle in $G \setminus W$ of order $\theta - |W|$.}

\textit{Proof.} Since $|W| < \theta$, the result is obvious when $\theta = 1$, and so we may assume that $\theta \ge 2$. Now $G \setminus W$ is obtained from $G/W$ by deleting edges with at most one end, and $\mathscr{T}'$ is obtained from $\mathscr{T}/W$ by repeating the operation of (8.4). The result follows. $\square$

\sectionhead{9. Laminar Separations}

We have seen in (5.2) that the tangles of large order are obstructions to the existence of tree-decompositions of small width. Our next result is a counterpart of this, that there is a tree-decomposition into pieces which correspond to the tangles.

Let $(A_1, B_1)$, $(A_2, B_2)$ be separations of a hypergraph $G$. We say these separations \textit{cross} unless either $A_1 \subseteq A_2$ and $B_2 \subseteq B_1$, or $A_1 \subseteq B_2$ and $A_2 \subseteq B_1$, or $B_1 \subseteq A_2$ and $B_2 \subseteq A_1$, or $B_1 \subseteq B_2$ and $A_2 \subseteq A_1$. A set of separations is \textit{laminar} if no two of its members cross.

Let $(T, \tau)$ be a tree-decomposition of a hypergraph $G$. For each $e \in E(T)$, let $T_1$, $T_2$ be the components of $T \setminus e$ and let
\begin{equation*}
G_i^e = \bigcup (\tau(t) : t \in V(T_i)) \qquad (i = 1, 2).
\end{equation*}
Then $(G_1^e, G_2^e)$ is a separation of $G$, and we call $(G_1^e, G_2^e)$ and $(G_2^e, G_1^e)$ the separations \textit{made by} $e$ (under the given tree-decomposition).

\stmt
(9.1) \textit{If $(T, \tau)$ is a tree-decomposition of $G$, then the set of all separations of $G$ made by edges of $T$ is laminar. Conversely, if $\{(A_i, B_i) : 1 \le i \le k\}$ is a laminar set of separations of $G$, there is a tree-decomposition $(T, \tau)$ of $G$ such that}

\hspace*{12pt}(i) \textit{for $1 \le i \le k$, $(A_i, B_i)$ is made by a unique edge of $T$}

\hspace*{12pt}(ii) \textit{for each edge $e$ of $T$, at least one of the two separations made by $e$ equals $(A_i, B_i)$ for some $i$ ($1 \le i \le k$).}

The proof is easy\del{ and is left to the reader}\fix{, and we give it now}.

\fix{For an edge $e$ of $T$ with ends $u, v$, let $G^e_u$ be the union of the subhypergraphs $\tau(t)$ over the vertices $t$ of the component of $T \setminus e$ that contains $u$. Thus $(G^e_u, G^e_v)$ and $(G^e_v, G^e_u)$ are the separations made by $e$. (They are separations, since $G^e_u \cup G^e_v = G$ by (i) of the definition of tree-decomposition, and no edge of $G$ is in both, by (ii).) If we replace one of two separations by its reverse, the four conditions in the definition of crossing are permuted. Hence two separations cross if and only if they cross after one or both of them are reversed.

For the first statement, let $e_1, e_2 \in E(T)$. If $e_1 = e_2$, the separations made by $e_1$ are $(A, B)$ and $(B, A)$ for some $A, B$, and two separations each equal to $(A, B)$ or $(B, A)$ do not cross: if they are equal the first condition holds, and otherwise the second holds. So we may assume that $e_1 \ne e_2$. Let $e_i$ have ends $u_i, v_i$ $(i = 1, 2)$. Since $e_2$ is an edge of $T \setminus e_1$, its ends are in the same component of $T \setminus e_1$, and by exchanging the names $u_1, v_1$ we may assume that this is the component that contains $v_1$. Similarly, by exchanging the names $u_2, v_2$, which keeps this property, we may assume that $u_1$ and $v_1$ are in the component of $T \setminus e_2$ that contains $u_2$. The component of $T \setminus e_1$ that contains $u_1$ does not contain $e_2$, since both ends of $e_2$ are in the other component, and so it is a connected subgraph of $T \setminus e_2$. It contains $u_1$, and so it is contained in the component of $T \setminus e_2$ that contains $u_2$. Hence $G^{e_1}_{u_1} \subseteq G^{e_2}_{u_2}$. Similarly the component of $T \setminus e_2$ that contains $v_2$ does not contain $e_1$, since both ends of $e_1$ are in the component that contains $u_2$. So it is a connected subgraph of $T \setminus e_1$, and it contains $v_2$, which is in the component of $T \setminus e_1$ that contains $v_1$. Hence $G^{e_2}_{v_2} \subseteq G^{e_1}_{v_1}$. So $(G^{e_1}_{u_1}, G^{e_1}_{v_1})$ and $(G^{e_2}_{u_2}, G^{e_2}_{v_2})$ satisfy the first condition and do not cross. The other separations made by $e_1$ and $e_2$ are their reverses, and so no two of the separations made by edges of $T$ cross.

For the converse, let $\mathscr{S}$ be the set of the separations $(A_i, B_i)$ and $(B_i, A_i)$ $(1 \le i \le k)$. It is laminar, by the remark on reversal. We prove by induction on $|\mathscr{S}|$ that if $\mathscr{S}$ is a laminar set of separations of $G$ that contains the reverse of each of its members, then there is a tree-decomposition $(T, \tau)$ of $G$ such that the separations made by edges of $T$ are exactly the members of $\mathscr{S}$, and each member of $\mathscr{S}$ is made by a unique edge of $T$. This proves (i) and (ii), since the two separations made by an edge of $T$ are the reverses of each other. If $\mathscr{S} = \emptyset$, let $T$ have one vertex $t$ and put $\tau(t) = G$. So we may assume that $\mathscr{S} \ne \emptyset$. For a separation $(A, B)$ of $G$ let
\begin{equation*}
s(A, B) = |V(A)| + |E(A)| + |V(G) - V(B)| + |E(G) - E(B)|,
\end{equation*}
and choose $(A_0, B_0) \in \mathscr{S}$ with $s(A_0, B_0)$ minimum. Let $\mathscr{S}'$ be obtained from $\mathscr{S}$ by removing $(A_0, B_0)$ and $(B_0, A_0)$. Then $\mathscr{S}'$ is laminar and contains the reverse of each of its members, and so from our inductive hypothesis there is a tree-decomposition $(T', \tau')$ of $G$ for $\mathscr{S}'$ as above. For an edge $e$ of $T'$ with ends $u, v$ we define $G'^e_u$ for $(T', \tau')$ as we defined $G^e_u$.

\hspace*{12pt}(1) \textit{For each $(C, D) \in \mathscr{S}'$, either $A_0 \subseteq C$ and $D \subseteq B_0$, or $A_0 \subseteq D$ and $C \subseteq B_0$.}

For $(A_0, B_0)$ and $(C, D)$ do not cross, and so one of the four conditions holds. The first two are the two outcomes of (1). Suppose that $B_0 \subseteq C$ and $D \subseteq A_0$. Then $(D, C) \in \mathscr{S}$, and $s(D, C) \le s(A_0, B_0)$ with equality only if $D = A_0$ and $C = B_0$, since $D \subseteq A_0$ and $B_0 \subseteq C$. By the choice of $(A_0, B_0)$ equality holds, and so $(C, D) = (B_0, A_0) \notin \mathscr{S}'$, a contradiction. Similarly, if $B_0 \subseteq D$ and $C \subseteq A_0$, then $(C, D) = (A_0, B_0)$, again a contradiction. This proves (1).

Call an ordered pair $(u, w)$ of adjacent vertices of $T'$ \textit{good} if $A_0 \subseteq G'^e_u$ and $G'^e_w \subseteq B_0$, where $e$ is the edge with ends $u, w$. Since $(G'^e_u, G'^e_w) \in \mathscr{S}'$, (1) implies that for every edge $e$ of $T'$ with ends $u, w$, at least one of $(u, w)$, $(w, u)$ is good.

\hspace*{12pt}(2) \textit{There is a vertex $u_0$ of $T'$ such that $(u_0, w)$ is good for every neighbour $w$ of $u_0$.}

For suppose not. Then each vertex $v$ of $T'$ has a neighbour $u$ such that $(v, u)$ is not good, and then $(u, v)$ is good. Starting from any vertex and repeating, we obtain an infinite walk $v_0 v_1 v_2 \dots$ of $T'$ such that for each $i$, $(v_i, v_{i+1})$ is not good and $(v_{i+1}, v_i)$ is good. Then $v_{i+2} \ne v_i$, since $(v_{i+1}, v_{i+2})$ is not good and $(v_{i+1}, v_i)$ is good. A walk in a tree that never goes back along the edge it has just used is a path, and so $T'$ has arbitrarily long paths, which is impossible. This proves (2).

\hspace*{12pt}(3) \textit{$\tau'(t) \subseteq B_0$ for every vertex $t \ne u_0$ of $T'$, and $A_0 \subseteq \tau'(u_0)$.}

Let $t \ne u_0$, and let $w$ be the neighbour of $u_0$ on the path of $T'$ from $u_0$ to $t$, and $e$ the edge with ends $u_0, w$. Then $t$ is in the component of $T' \setminus e$ that contains $w$, and so $\tau'(t) \subseteq G'^e_w \subseteq B_0$ by (2). For the second statement, let $x \in V(A_0)$. Then $x \in V(\tau'(t))$ for some $t$. If $t \ne u_0$, define $w$ and $e$ as before. Since $A_0 \subseteq G'^e_{u_0}$ by (2), $x \in V(\tau'(t_1))$ for some vertex $t_1$ of the component of $T' \setminus e$ that contains $u_0$. Then $u_0$ is on the path of $T'$ between $t_1$ and $t$, and so $x \in V(\tau'(u_0))$. Now let $f \in E(A_0)$. Then $f \in E(\tau'(t))$ for some $t$, and $t = u_0$, since otherwise $f \in E(B_0)$ by the first statement, which is impossible since $E(A_0 \cap B_0) = \emptyset$. This proves (3).

Now let $T$ be obtained from $T'$ by adding a new vertex $t_0$ and an edge $e_0$ joining $t_0$ and $u_0$. Let $\tau(t_0) = A_0$, $\tau(u_0) = \tau'(u_0) \cap B_0$, and $\tau(t) = \tau'(t)$ for every other vertex $t$. Since $\tau'(u_0) \subseteq A_0 \cup B_0$ and $A_0 \subseteq \tau'(u_0)$ by (3),
\begin{equation*}
\tau(t_0) \cup \tau(u_0) = A_0 \cup (\tau'(u_0) \cap B_0) = \tau'(u_0).
\end{equation*}
We check that $(T, \tau)$ is a tree-decomposition of $G$. First, $T$ is a tree, and (i) holds by (i) for $(T', \tau')$ and the equation above. For (ii), $\tau(t_0) = A_0$ has no edge in common with $\tau(u_0)$ or with any $\tau(t) = \tau'(t)$ $(t \ne t_0, u_0)$, since these are subhypergraphs of $B_0$ by (3) and $E(A_0 \cap B_0) = \emptyset$. The other pairs satisfy (ii) since $\tau(t) \subseteq \tau'(t)$ for every vertex $t$ of $T'$. For (iii), let $t'$ be on the path of $T$ between $t$ and $t''$, where $t' \ne t, t''$ (otherwise (iii) is trivial). Then $t' \ne t_0$, since $t_0$ has valency 1. If $t_0 \notin \{t, t''\}$, then $\tau(t) \cap \tau(t'') \subseteq \tau'(t')$ by (iii) for $(T', \tau')$. If $t_0 \in \{t, t''\}$, say $t = t_0$, then $t'$ is on the path of $T'$ between $u_0$ and $t''$. Since $A_0 \subseteq \tau'(u_0)$, it follows that $\tau(t_0) \cap \tau(t'') \subseteq \tau'(u_0) \cap \tau'(t'') \subseteq \tau'(t')$. So in each case $\tau(t) \cap \tau(t'') \subseteq \tau'(t')$. This proves (iii) when $t' \ne u_0$. When $t' = u_0$, neither $t$ nor $t''$ is $u_0$, and one of them is not $t_0$. That one is a vertex $t_1 \ne u_0$ of $T'$, and $\tau(t_1) \subseteq B_0$ by (3). Hence $\tau(t) \cap \tau(t'') \subseteq \tau'(u_0) \cap B_0 = \tau(u_0)$. So $(T, \tau)$ is a tree-decomposition of $G$.

The separations made by $e_0$ are $(A_0, B_0)$ and $(B_0, A_0)$. For the component of $T \setminus e_0$ that contains $t_0$ has only the vertex $t_0$, and $\tau(t_0) = A_0$. The union of $\tau(t)$ over the other vertices is a subhypergraph of $B_0$ by (3). It also contains every vertex and edge of $B_0$, since each of them is in some $\tau'(t)$, and if $t = u_0$ it is in $\tau'(u_0) \cap B_0 = \tau(u_0)$. So this union is $B_0$. Now let $e$ be an edge of $T'$. Then $t_0$ and $u_0$ are in the same component of $T \setminus e$, and $\tau(t_0) \cup \tau(u_0) = \tau'(u_0)$. So $e$ makes the same separations under $(T, \tau)$ as under $(T', \tau')$. Hence the separations made by edges of $T$ are the members of $\mathscr{S}' \cup \{(A_0, B_0), (B_0, A_0)\} = \mathscr{S}$. Each member of $\mathscr{S}'$ is made by a unique edge of $T'$, and not by $e_0$. The separations $(A_0, B_0)$ and $(B_0, A_0)$ are made by $e_0$, and by no edge of $T'$, since they are not in $\mathscr{S}'$. This completes the induction, and so proves (9.1).}\fixnote{fix:9.1-proof-added}{The paper leaves this proof to the reader. The argument above is the one used in the formalization.}

We wish to arrange a ``tie-breaking'' mechanism so that no two distinct separations are counted as having the same order (except for reversal). A \textit{tie-breaker} $\lambda$ in a hypergraph $G$ is a function from the set of all separations of $G$ into some linearly ordered set $(\Lambda, <)$, satisfying certain axioms given below. For each separation $(A, B)$, $\lambda(A, B)$ is called the $\lambda$-\textit{order} of $(A, B)$,
